\section{The \texttt{Bragg\_crystal\_bent\_BC} McXtrace Component}
Bent, perfect crystal with common cubic structures (diamond, fcc, or bcc, and others if symmetry form factor multipliers provided explicitly)

\subsection*{Identification}
\begin{itemize}
  \item \textbf{Author:} Marcus H Mendenhall, NIST \textless{}marcus.mendenhall@nist.gov\textgreater{}
  \item \textbf{Origin:} NIST
  \item \textbf{Date:} Dec 2016
\end{itemize}

\subsection*{Description}
Bragg\_crystal\_bent\_BC.com is intended to supercede Bragg\_Crystal\_bent.comp For details see: The optics of focusing bent-crystal monochromators on X-ray powder diffractometers with application to lattice parameter determination and microstructure analysis, Marcus H. Mendenhall,* David Black and James P. Cline, J. Appl. Cryst. (2019). 52, https://doi.org/10.1107/S1600576719010951

Reads atomic formfactors from a data input file. The Bragg\_Crystal code reflects ray in an ideal geometry, does not include surface imperfections or mosaicity

The crystal code reflects ray in an ideal geometry, i.e. does not include surface imperfections or mosaicity. The crystal planes from which the reflection is made lies in the X-Z plane on the unbent crystal rotated by an angle alpha about the Y axis with respect to the crystal surface.

The crystal itself is set in the X-Z plane positioned such that the long axis of the crystal surface coincides with the Z-axis, withs normal pointing in the poisitivce Y-direction.

The asummetry angle alpha is defined so that positive alpha reduces the Bragg angle to the plane i.e. alpha=Thetain grazes the planes. if alpha!=0, one should restrict to rays which have small kx values, since otherwise the alpha rotation is not around the diffraction axis.

The mirror is positioned such that the a-axis of the mirror ellipsoid is on the z-axis, the b-axis is along the y-axis and the c is along the x-axis. The reference point of the mirror is the ellipsoid centre, offset by one half-axis along the y-axis. (See the component manual for a drawing).

Notation follows Tadashi Matsushita and Hiro-O Hashizume, X-RAY MONOCHROMATORS. Handbook on Synchrotron Radiation,North-Holland Publishing Company, 1:263–274, 1983.

NOTE: elliptical coordinate code and documentation taken from Mirror\_elliptic.comp distributed in McXtrace v1.2 written by: Erik Knudsen. However, the coordinates are rotated to be consistent with Perfect\_Crystal.comp and NIST\_Perfect\_Crystal.comp Idealized elliptic mirror with surface ellipse and lattice ellipses independent, to allow construction of Johansson optics, for example.

Non-copyright notice: Contributed by the National Institute of Standards and Technology; not subject to copyright in the United States. This is not an official contribution, in that the results are in no way certified by NIST.

Example: Bragg\_crystal\_bent\_BC( length=0.05, width=0.02, V=160.1826, h=1, k=1, l=1, alpha=0)

\subsection*{Input parameters}
Parameters in \textbf{boldface} are required; the others are optional.

\begin{longtable}{p{0.22\textwidth}p{0.12\textwidth}p{0.46\textwidth}p{0.14\textwidth}}
\toprule
\textbf{Name} & \textbf{Unit} & \textbf{Description} & \textbf{Default} \\
\midrule
\endhead
x\_a & m & 1st short half axis (along x). Commonly set to zero, which really implies infinite value, so crystal is an elliptic cylinder. & 0 \\
y\_b & m & 2nd short half axis (along y), which is also the presumed near-normal direction, reflection near the y-z plane. & 1.0 \\
z\_c & m & Long half axis (along z). Commonly a=0. b=c, which creates a circular cylindrical surface. & 1.0 \\
lattice\_x\_a & m & Curvature matrix component around a for underlying lattice, for bent/ground/rebent crystals & 0 \\
lattice\_y\_b & m & Curvature matrix component around b for underlying lattice, for bent/ground/rebent crystals & 1.0 \\
lattice\_z\_c & m & curvature matrix component around c for underlying lattice, for bent/ground/rebent crystals THERE HAS BEEN NO TESTING for the case in which lattice\_x\_a != x\_a. & 1.0 \\
length & m & z depth (length) of the crystal. & 0.05 \\
width & m & x width of the crystal. & 0.02 \\
V & \AA{}$^{3}$ & unit cell volume & 160.1826 \\
form\_factors &  &  & "FormFactors.txt" \\
material &   & Si, Ge (maybe also GaAs?) & "Si.txt" \\
alpha & rad & Asymmetry angle (alpha=0 for symmetric reflection, i.e. the Bragg planes are parallel to the crystal surface) & 0.0 \\
R0 &   & Reflectivity. Overrides the computed Darwin reflectivity. Probably only useful for debugging. & 0 \\
debye\_waller\_B & \AA{}$^{2}$ & Debye-Waller temperature factor, M=B*(sin(theta)/lambda)\textasciicircum{}2*(2/3), default=silicon at room temp. & 0.4632 \\
crystal\_type &   & 1 =\textgreater{} Mx\_crystal\_explicit: provide explicit real and imaginary form factor multipliers structure\_factor\_scale\_r, structure\_factor\_scale\_i; 2 =\textgreater{} Mx\_crystal\_diamond: diamond; 3 =\textgreater{} Mx\_crystal\_fcc: fcc; 4 =\textgreater{} Mx\_crystal\_fcc: bcc & 1 \\
h &   & Miller index of reflection & 1 \\
k &   & Miller index of reflection & 1 \\
l &   & Miller index of reflection & 1 \\
structure\_factor\_scale\_r &  &  & 0.0 \\
structure\_factor\_scale\_i &  &  & 0.0 \\
verbose &   & if non-zero: Output more information (warnings and messages) to the console. & 0 \\
\bottomrule
\end{longtable}

\subsection*{Links}
\begin{itemize}
  \item Component source code found in file \texttt{Bragg\_crystal\_bent\_BC.comp}.
  \item material datafile obtained from http://physics.nist.gov/cgi-bin/ffast/ffast.pl
\end{itemize}
\IfFileExists{contrib/Bragg_crystal_bent_BC_static.tex}{\input{contrib/Bragg_crystal_bent_BC_static.tex}}{}